\pagenumbering{gobble}

\chapter*{Il percorso}
È passato tanto tempo da quando è iniziato questo percorso: tanto tempo in cui ho dovuto imparare a provvedere da solo a cosa mangiare, come organizzare il proprio tempo, quali momenti dedicare allo studio, quali al lavoro, quali alle persone che si amano e quali a se stessi.
È vero che quando non hai più qualcosa che davi per scontato, ti rendi conto di quanto, forse, scontato non lo era.

Un enorme sacrificio, seguito da un investimento di tempo, denaro ed energie, tutto per inseguire la promessa di un futuro dignitoso.
Numerosi anni di problemi da risolvere, decisioni da prendere - quasi mai semplici - un continuo compromesso tra quello che bisogna fare e quello che si vorrebbe fare.

La cosa più importante che ho imparato è stata prendere decisioni.
La seconda, una volta imparata la prima, è stata gestire gli imprevisti: e quindi anche saper rimettere in discussione le decisioni già prese.

Durante questo lungo viaggio, durato cinque anni, ci sono state tantissime difficoltà, tutte risolte (sperabilmente) nel migliore dei modi.

Nonostante l'apparente sicurezza che sembra scaturire dalle righe sopra, il percorso non è stato sempre così definito, e ha subito qualche ripensamento lungo la strada.
Anche ora sono contento di dove sono, ma non è detto che rimarrà tutto così anche in futuro.

In particolare, tra dicembre 2025 e marzo 2026, quindi a ridosso di Natale, avevo sul tavolo diverse proposte da docenti e aziende: collaborazioni, dottorati, contratti di lavoro.
Qualcuno direbbe \"Fantastico\", altri sarebbero entusiasti, ad altri ancora non importerebbe nulla, e continuerebbero a seguire il proprio piano di vita, scritto su un foglietto di carta all'età di 18 anni.

Io invece sono stato travolto da un'ondata di panico totale, come se dovessi scegliere per tutta la vita cosa fare - ed effettivamente era così. \textbf{MA}, papà e mamma mi hanno detto che nella vita si può sempre cambiare. E chi meglio di loro sa cosa significa \textit{cambiare}: in un modo o nell'altro si sono sempre dati da fare, cambiando come lavorano e cosa fanno per adattarsi alle esigenze del momento. Non è scontato, non è facile, e soprattutto mi hanno insegnato - mai a parole, ma con i fatti - che la vita non è un contratto di assicurazione.

Quindi, se loro possono cambiare il loro modo di vivere, perché non posso farlo anche io?

Beh, non posso farlo perché\ldots
\textit{e se poi non mi piace?},
\textit{e se poi mi pento?},
\textit{e se poi sono precario fino a 700 (settecento) anni?},
\textit{e se poi Elisa non mi vuole più?},
\textit{e se finisco sotto un ponte?},
\textit{e se va tutto male?}

Il fatto che a me non piaccia fare qualcosa, o che poi sia precario, o che me ne penta, è una cosa che può capitare, purtroppo - ma almeno poi so cosa \textbf{NON} voglio fare.
E in ogni situazione \textit{scomoda} in cui mi sono messo (o mi hanno messo) me la sono sempre cavata: sbagliare è umano.

La parte qui sopra è stata scritta prima di cambiare lavoro e prima di prendere le decisioni più importanti, MOLTO prima dell'inizio della tesi e della fine di tutti gli esami.
È stato divertente rileggere quello che il Daniele del passato mi aveva lasciato scritto.
Però non si vive di ricordi, Elisa me lo ripete sempre - ma che ci vuoi fare, mi piace ricordare i tempi andati: d'altronde siamo quello che siamo grazie ai ricordi, no?

Adesso basta così, però.
Torniamo ai giorni nostri.

\section*{Ringraziamenti}

Ringraziare tutti è difficile (non \texttt{NP}-difficile - RIDETE).

Battutoni a parte, ringrazio il mio relatore, il professor Ferretti, per avermi seguito nella realizzazione della tesi, per avermi fornito tante idee su come muovermi nella direzione che volevo e per aver risposto alle mie domande via call - domande per cui non sarebbe servita una call, e probabilmente nemmeno una mail, ma la pazienza è una virtù dei forti.
Ringrazio anche il professor Bedogni per avermi seguito nella fase iniziale del lavoro.

Un grazie enorme va ai miei genitori, mamma e papà, che ci sono sempre stati, con una pazienza variabile a seconda dei giorni della settimana.
Li ringrazio anche perché hanno fatto uno sforzo ENORME nel vedere il proprio figlio andare via - un po' come un bambino che guarda il proprio palloncino volare sempre più in alto, e prova comunque a rincorrerlo anche se ormai è fuori portata.
Da mamma probabilmente ho ereditato la forza di non arrendermi mai, da papà una proverbiale pazienza quasi infinita.

Grazie anche ad Alessio e Manuel, i miei fratelli: nonostante litighiamo (quasi) sempre, hanno fatto un sacco di strada per venirmi a trovare, e sono contento della loro presenza anche se spesso i litigi nascono per motivi di poco conto.
Voglio anche che sappiano che, per quanto litighino con mamma e papà, in realtà tengono a tutti noi - quindi ogni tanto ci vuole pazienza con i nostri vecchietti, anche quando \textbf{NON ASCOLTANO}.

Un grazie va anche ai miei nonni: Angela, Antonietta, Vittorio, che mi supportano chiamandomi, e Domenico, che mi segue da un po' più lontano; e ai miei zii, per le chiacchiere e il supporto (sopportato più che altro) che mi hanno dato in questi anni.

Vorrei poi ringraziare Elisa, la mia ragazza, che mi ama tanto e mi supporta tantissimo.
Non mi chiede molto, ma su alcune cose è abbastanza esigente: vuole che io mi prenda cura di me - sarebbe anche ora, direbbe qualcuno, visto quanto posso essere sciatto certe volte.
Elisa, oltre ad amarmi tanto, ha il difficilissimo compito di non farmi rispondere a caso alle persone: compito arduo, perché spesso capita che io non sia troppo ragionevole, ma sa come prendermi, e infatti riesce molto bene nel suo difficilissimo compito.
Ringrazio anche Lorenzo e tutta la loro famiglia per sopportare i miei discorsi la domenica a pranzo.

Un grazie anche a Giovanni, un grande amico che dai tempi delle medie mi accompagna nelle mie avventure nei server di Minecraft lagganti e non solo: probabilmente dovrei essere più presente, anzi sicuramente dovrei esserlo, anche se senza l'onere della laurea sarà molto più facile giocare e spettegolare come due signore anziane di paese.

Un pensiero va ai miei amici di Campobasso (CB), che hanno fatto tanti chilometri per vedere\ldots me: Danilo, Ilaria, Maria e Rosanna.
E a chi avrebbe voluto esserci ma non è potuto venire per vari motivi: Flavia, Franco, Giò Petti, Licia, Luca, Mattia, Marco e Samuele.

Grazie ai miei amici informatici, con cui ho sperimentato il brivido di parlare male del progetto di Sistemi Operativi: Alessandro R., Alessandro T., Camilla, Enrico, Francesco, Giulia, Leonardo, Lorenzo, Mattia (che come diceva Renato Carosone vorrebbe fare l'americano), Pietro e Simone.
E al gruppo di studio più o meno nuovo: Beatrice, Elena, Federica, Francesca, Gloria, Liam, Mariagrazia, Roberta, Rossella (tanti auguri) e Tommaso.

Grazie anche ai nuovi e ai vecchi coinquilini (che se sono qui non erano SOLO coinquilini, ma anche amichetti) e alle altre persone incontrate durante questo luuungo periodo: Corrado, Carola, Giulia e Salvatore.

Infine, si ringraziano TUTTE le altre persone non citate, i vecchi colleghi e i nuovi, e si ringraziano altresì tutti gli eventi non piacevoli che mi hanno insegnato cosa non voglio nella vita.

Io non so cosa mi riserverà il futuro.
L'unica cosa che so per certo è che sarà divertente, come lo è stato fino ad ora: un continuo divertimento, nonostante tutto.